
\section{SummEval Reproduction Study}\label{sec:6-chapter-summeval-repro}

We begin our reproduction study using the code-base released by the original authors.\footnote{\url{https://github.com/Yale-LILY/SummEval}} The file with the original metric scores used to calculate the system-level correlations reported in the paper is no longer available online. However, we reached out to the NLP community and a fellow researcher shared a copy of the file they had locally. We verified the accuracy of this file by using it to reproduce the scores reported in the original paper. We use the SummEval evaluation package and released script to re-calculate the metric scores and system-level correlations. For each metric, we use the parameter settings released in the SummEval code-base.

We report the system-level Kendall-Tau correlation of the evaluation metrics with the expert annotations, for both the original study and our replication study in~\Cref{tab:comparison}. As we can see from the results, we were not able to perfectly reproduce any of the metrics. While some amount of difference can be expected due to numerical instability, many of the differences in our reproduction study are too large to attribute to simple instability. For example, the original study reports a correlation of $0.074$ for ROUGE-L and Coherence. However, in our study we find a correlation of $0.417$, over a $0.3$ difference in correlation. Similarly, the BLEU-Coherence correlation shows an increase of $0.265$ in our reproduction study. In general, we find that ROUGE, BLEU, and n-gram based measures have the largest differences in our reproduction study. We hypothesize that this is because ROUGE and BLEU rely on strict n-gram matching, which is highly sensitive to changes in tokenization settings. The remaining metrics have much lower differences in our study. We also note that we were unable to reproduce four metrics reported in the original study ($\mathrm{S}^3$, ROUGE-we, SMS, and MoverScore), due to broken links to files required to run the metrics. We note that all of the metrics we were unable to reproduce are in the Embedding/neural metric category, pointing to a general problem of reproducibility with this style of metric.


It is important that we design metrics that allow for reproducible evaluation, both in design and practice. Regarding design, we show that evaluation metrics that rely on strict token overlap are difficult to reproduce due to their sensitivity to tokenization settings. Additionally, metrics that rely on pre-trained models or embeddings need to be easily accessible or easily reproducible, otherwise they create a barrier to use. Regarding practice, as a community we have made great strides towards making code release a norm for publication. However, we need additional research studying the instability of evaluation metrics.


% -----------------------------------------------------------------------
% Comparison Table: All Metrics on SummEval
% Each metric appears once at the maximum nrefs available.
% -----------------------------------------------------------------------
\begin{sidewaystable}[h]
% \begin{table*}[h]
\centering
\scriptsize
% \setlength{\tabcolsep}{6pt}
\renewcommand{\arraystretch}{0.9}
\begin{tabular}{l r
    ccc
    ccc
    ccc
    ccc}
\toprule
& & \multicolumn{3}{c}{\textbf{Coherence}}
& \multicolumn{3}{c}{\textbf{Consistency}}
& \multicolumn{3}{c}{\textbf{Fluency}}
& \multicolumn{3}{c}{\textbf{Relevance}} \\
\cmidrule(lr){3-5}\cmidrule(lr){6-8}\cmidrule(lr){9-11}\cmidrule(lr){12-14}
\textbf{Metric} & \textbf{$n$}
  & \textit{Orig.} & \textit{Repl.} & \textit{$\Delta$}
  & \textit{Orig.} & \textit{Repl.} & \textit{$\Delta$}
  & \textit{Orig.} & \textit{Repl.} & \textit{$\Delta$}
  & \textit{Orig.} & \textit{Repl.} & \textit{$\Delta$} \\
\midrule
\multicolumn{14}{l}{\textit{String-match / overlap}} \\[2pt] 
ROUGE-1
  & 11
  & 0.250 & 0.367 & \dpos{0.117}
  & 0.529 & 0.567 & \dpos{0.038}
  & 0.524 & 0.550 & \dpos{0.026}
  & 0.412 & 0.600 & \dpos{0.188} \\

ROUGE-2
  & 11
  & 0.162 & 0.233 & \dpos{0.071}
  & 0.588 & 0.600 & \dpos{0.012}
  & 0.480 & 0.483 & \dpos{0.003}
  & 0.294 & 0.433 & \dpos{0.139} \\

ROUGE-3
  & 11
  & 0.221 & 0.300 & \dpos{0.079}
  & 0.706 & \textbf{0.733} & \dpos{0.027}
  & 0.509 & 0.517 & \dpos{0.008}
  & 0.353 & 0.500 & \dpos{0.147} \\

ROUGE-4
  & 11
  & 0.309 & 0.383 & \dpos{0.074}
  & 0.588 & 0.617 & \dpos{0.029}
  & 0.553 & 0.533 & \dneg{0.020}
  & 0.412 & 0.550 & \dpos{0.138} \\

ROUGE-L
  & 11
  & 0.073 & 0.417 & \dpos{0.344}
  & 0.147 & 0.283 & \dpos{0.136}
  & 0.258 & 0.500 & \dpos{0.242}
  & 0.235 & 0.550 & \dpos{0.315} \\

ROUGE-su*
  & 11
  & 0.191 & 0.300 & \dpos{0.109}
  & 0.294 & 0.300 & \dpos{0.006}
  & 0.435 & 0.450 & \dpos{0.015}
  & 0.324 & 0.500 & \dpos{0.176} \\

ROUGE-w
  & 11
  & 0.000 & 0.350 & \dpos{0.350}
  & 0.397 & 0.683 & \dpos{0.286}
  & 0.376 & 0.533 & \dpos{0.157}
  & 0.162 & 0.550 & \dpos{0.388} \\


BLEU
  & 11
  & 0.118 & 0.383 & \dpos{0.265}
  & 0.073 & $-$0.050 & \dneg{0.123}
  & 0.332 & 0.233 & \dneg{0.099}
  & 0.221 & 0.283 & \dpos{0.062} \\

CHRF
  & 11
  & 0.397 & 0.350 & \dneg{0.047}
  & 0.529 & 0.617 & \dpos{0.088}
  & 0.465 & 0.567 & \dpos{0.102}
  & 0.588 & 0.550 & \dneg{0.038} \\

CIDEr
  & 11
  & 0.118 & 0.067 & \dneg{0.051}
  & $-$0.191 & $-$0.167 & \dpos{0.024}
  & $-$0.022 & 0.050 & \dpos{0.072}
  & 0.191 & 0.067 & \dneg{0.124} \\

METEOR
  & 11
  & 0.235 & 0.283 & \dpos{0.048}
  & 0.632 & 0.650 & \dpos{0.018}
  & 0.613 & 0.600 & \dneg{0.013}
  & 0.426 & 0.550 & \dpos{0.124} \\
\midrule
\multicolumn{14}{l}{\textit{Embedding / neural}} \\[2pt]

ROUGE-we-1$^{\dagger}$
  & 11
  & 0.265 & --- & \dna
  & 0.456 & --- & \dna
  & 0.509 & --- & \dna
  & 0.426 & --- & \dna \\

ROUGE-we-2$^{\dagger}$
  & 11
  & $-$0.015 & --- & \dna
  & 0.500 & --- & \dna
  & 0.303 & --- & \dna
  & 0.118 & --- & \dna \\

ROUGE-we-3$^{\dagger}$
  & 11
  & 0.029 & --- & \dna
  & 0.368 & --- & \dna
  & 0.303 & --- & \dna
  & 0.191 & --- & \dna \\

$S^3$-pyr$^{\dagger}$
  & 11
  & $-$0.029 & --- & \dna
  & 0.515 & --- & \dna
  & 0.317 & --- & \dna
  & 0.132 & --- & \dna \\

$S^3$-resp$^{\dagger}$
  & 11
  & $-$0.015 & --- & \dna
  & 0.500 & --- & \dna
  & 0.332 & --- & \dna
  & 0.147 & --- & \dna \\


MoverScore$^{\dagger}$
  & 11
  & 0.191 & --- & \dna
  & $-$0.029 & --- & \dna
  & 0.258 & --- & \dna
  & 0.294 & --- & \dna \\

SMS$^{\dagger}$
  & 11
  & 0.162 & --- & \dna
  & 0.559 & --- & \dna
  & 0.362 & --- & \dna
  & 0.235 & --- & \dna \\

BERTScore-P
  & 11
  & 0.059 & $-$0.017 & \dneg{0.076}
  & $-$0.191 & $-$0.217 & \dneg{0.026}
  & 0.007 & 0.033 & \dpos{0.026}
  & 0.162 & 0.083 & \dneg{0.079} \\

BERTScore-R
  & 11
  & 0.147 & 0.200 & \dpos{0.053}
  & 0.662 & 0.667 & \dpos{0.005}
  & 0.494 & 0.483 & \dneg{0.011}
  & 0.309 & 0.433 & \dpos{0.124} \\

BERTScore-F
  & 11
  & 0.206 & 0.133 & \dneg{0.073}
  & 0.044 & 0.033 & \dneg{0.011}
  & 0.243 & 0.283 & \dpos{0.040}
  & 0.426 & 0.400 & \dneg{0.026} \\


SummaQA
  & 1
  & 0.118 & 0.233 & \dpos{0.115}
  & 0.603 & 0.700 & \dpos{0.097}
  & 0.406 & 0.494 & \dpos{0.088}
  & 0.221 & 0.433 & \dpos{0.212} \\

SuPERT
  & 1
  & 0.103 & 0.200 & \dpos{0.097}
  & 0.588 & 0.600 & \dpos{0.012}
  & 0.421 & 0.383 & \dneg{0.038}
  & 0.235 & 0.400 & \dpos{0.165} \\

BLANC
  & 0
  & 0.073 & 0.183 & \dpos{0.110}
  & 0.559 & 0.617 & \dpos{0.058}
  & 0.362 & 0.410 & \dpos{0.048}
  & 0.265 & 0.383 & \dpos{0.118} \\



\midrule
\multicolumn{14}{l}{\textit{Structural statistics}} \\[2pt]

Length
  & 0
  & $-$0.029 & 0.000 & \dpos{0.029}
  & 0.426 & 0.467 & \dpos{0.041}
  & 0.258 & 0.226 & \dneg{0.032}
  & 0.162 & 0.233 & \dpos{0.071} \\

Novel unigram
  & 0
  & 0.147 & 0.150 & \dpos{0.003}
  & $-$0.221 & $-$0.050 & \dpos{0.171}
  & $-$0.140 & 0.126 & \dpos{0.266}
  & 0.103 & 0.083 & \dneg{0.020} \\

Novel bi-gram
  & 0
  & 0.029 & $-$0.100 & \dneg{0.129}
  & $-$0.544 & $-$0.467 & \dpos{0.077}
  & $-$0.347 & $-$0.193 & \dpos{0.154}
  & $-$0.103 & $-$0.200 & \dneg{0.097} \\

Novel tri-gram
  & 0
  & 0.029 & $-$0.100 & \dneg{0.129}
  & $-$0.574 & $-$0.533 & \dpos{0.041}
  & $-$0.347 & $-$0.259 & \dpos{0.088}
  & $-$0.132 & $-$0.267 & \dneg{0.135} \\

Repeated unigram
  & 0
  & $-$0.382 & $-$0.333 & \dpos{0.049}
  & 0.103 & 0.067 & \dneg{0.036}
  & $-$0.066 & $-$0.142 & \dneg{0.076}
  & $-$0.368 & $-$0.300 & \dpos{0.068} \\

Repeated bi-gram
  & 0
  & $-$0.382 & $-$0.350 & \dpos{0.032}
  & $-$0.015 & $-$0.050 & \dneg{0.035}
  & $-$0.243 & $-$0.293 & \dneg{0.050}
  & $-$0.456 & $-$0.417 & \dpos{0.039} \\

Repeated tri-gram
  & 0
  & $-$0.221 & $-$0.167 & \dpos{0.054}
  & 0.147 & 0.133 & \dneg{0.014}
  & $-$0.022 & $-$0.075 & \dneg{0.053}
  & $-$0.265 & $-$0.200 & \dpos{0.065} \\

Stats-coverage
  & 0
  & $-$0.132 & $-$0.067 & \dpos{0.065}
  & 0.353 & 0.100 & \dneg{0.253}
  & 0.155 & $-$0.042 & \dneg{0.197}
  & $-$0.029 & 0.033 & \dpos{0.062} \\

Stats-compression
  & 0
  & 0.118 & 0.067 & \dneg{0.051}
  & $-$0.426 & $-$0.433 & \dneg{0.007}
  & $-$0.229 & $-$0.193 & \dpos{0.036}
  & $-$0.015 & $-$0.133 & \dneg{0.118} \\

Stats-density
  & 0
  & 0.162 & 0.200 & \dpos{0.038}
  & 0.647 & 0.633 & \dneg{0.014}
  & 0.391 & 0.360 & \dneg{0.031}
  & 0.294 & 0.367 & \dpos{0.073} \\

\bottomrule
% \multicolumn{14}{l}{%
%   \footnotesize
%   $n$ = number of references used (0 = reference-free). \\
%   $^{\dagger}$ Not reproduced; original values only. \\
%   $^{\ddagger}$ Not in original study; replication values only. \\
%   For SuPERT, original used $n{=}1$ and replication used $n{=}11$ (maximum available per study).
% } \\
\end{tabular}
\caption[SummEval reproducibility study results.]{%
    \footnotesize
  \textbf{All Metrics on SummEval} (Gold references, maximum nrefs available per metric).
  Comparison between the \emph{original} study and our \emph{replication}.
  $\Delta$ = Replication $-$ Original;
  \textcolor{deltapos}{green} $=$ replication higher,
  \textcolor{deltaneg}{red} $=$ replication lower.
  $^{\dagger}$ Metric not reproduced; original values only.
}
\label{tab:comparison}
% \end{table*}
\end{sidewaystable}

\clearpage

% -----------------------------------------------------------------------
% TABLE: Traditional Metrics on SummEval (non-LLM), best nrefs per metric
% -----------------------------------------------------------------------
\begin{table*}[h!]
\scriptsize
\centering
\setlength{\tabcolsep}{1pt}
\begin{subtable}[t]{0.49\textwidth}
\centering
\begin{tabular}{lllcccc}
\toprule
\textbf{Metric} & \textbf{Ref.} & \textbf{nrefs}
  & \textbf{Coh.} & \textbf{Con.} & \textbf{Flu.} & \textbf{Rel.} \\
\midrule
\rowcolor{refbased}
ROUGE-1 & Gold & 6 & 0.450 & 0.450 & 0.561 & 0.650 \\
\rowcolor{refbased}
ROUGE-2 & Gold & 6 & 0.333 & 0.500 & 0.577 & 0.533 \\
\rowcolor{refbased}
ROUGE-3 & Gold & 6 & 0.383 & 0.683 & 0.644 & 0.550 \\
\rowcolor{refbased}
ROUGE-4 & Gold & 6 & 0.500 & 0.567 & \textbf{0.711} & \textbf{0.700} \\
\rowcolor{refbased}
ROUGE-L & Gold & 6 & \textbf{0.550} & 0.317 & 0.628 & 0.683 \\
\rowcolor{refbased}
ROUGE-su* & Gold & 6 & 0.500 & 0.333 & 0.577 & 0.633 \\
\rowcolor{refbased}
ROUGE-w & Gold & 6 & 0.417 & 0.617 & 0.561 & 0.617 \\
\rowcolor{refbased}
BertScore-p & Gold & 11 & -0.017 & -0.217 & 0.033 & 0.083 \\
\rowcolor{refbased}
BertScore-r & Gold & 6 & 0.317 & 0.683 & 0.544 & 0.483 \\
\rowcolor{refbased}
BertScore-f & Gold & 11 & 0.133 & 0.033 & 0.283 & 0.400 \\
\rowcolor{refbased}
BLEU & Gold & 6 & 0.433 & -0.067 & 0.243 & 0.300 \\
\rowcolor{refbased}
CHRF & Gold & 11 & 0.350 & 0.617 & 0.567 & 0.550 \\
\rowcolor{refbased}
CIDEr & Gold & 1 & 0.083 & -0.083 & 0.126 & 0.150 \\
\rowcolor{refbased}
METEOR & Gold & 6 & 0.300 & 0.667 & 0.628 & 0.533 \\
\rowcolor{refbased}
SummaQA & Gold & 1 & 0.233 & \textbf{0.700} & 0.494 & 0.433 \\
\rowcolor{refbased}
SuPERT & Gold & 11 & 0.200 & 0.600 & 0.393 & 0.400 \\
\bottomrule
\end{tabular}
\caption{Reference-based: Gold (Standard) references.}
\label{tab:traditional_metrics-refbased}
\end{subtable}%
\hfill
\begin{subtable}[t]{0.49\textwidth}
\centering
\begin{tabular}{lllcccc}
\toprule
\textbf{Metric} & \textbf{Ref.} & \textbf{nrefs}
  & \textbf{Coh.} & \textbf{Con.} & \textbf{Flu.} & \textbf{Rel.} \\
\midrule
\rowcolor{refadj}
ROUGE-1 & Abstractive & 6 & 0.427 & 0.471 & 0.583 & 0.559 \\
\rowcolor{refadj}
ROUGE-2 & Abstractive & 6 & 0.324 & 0.515 & 0.598 & 0.456 \\
\rowcolor{refadj}
ROUGE-3 & Extractive & 6 & 0.353 & \textbf{0.691} & 0.647 & 0.456 \\
\rowcolor{refadj}
ROUGE-4 & Abstractive & 6 & 0.456 & 0.588 & \textbf{0.701} & \textbf{0.588} \\
\rowcolor{refadj}
ROUGE-L & Abstractive & 6 & \textbf{0.515} & 0.353 & 0.613 & \textbf{0.588} \\
\rowcolor{refadj}
ROUGE-su* & Abstractive & 6 & 0.456 & 0.353 & 0.554 & 0.529 \\
\rowcolor{refadj}
ROUGE-w & Abstractive & 6 & 0.353 & 0.632 & 0.598 & 0.485 \\
\rowcolor{refadj}
BertScore-p & Abstractive & 11 & 0.059 & -0.191 & 0.007 & 0.162 \\
\rowcolor{refadj}
BertScore-r & Extractive & 6 & 0.309 & 0.677 & 0.574 & 0.412 \\
\rowcolor{refadj}
BertScore-f & Abstractive & 11 & 0.206 & 0.044 & 0.244 & 0.427 \\
\rowcolor{refadj}
BLEU & Abstractive & 6 & 0.471 & -0.074 & 0.185 & 0.368 \\
\rowcolor{refadj}
CHRF & Abstractive & 6 & 0.397 & 0.529 & 0.465 & \textbf{0.588} \\
\rowcolor{refadj}
CIDEr & Abstractive & 1 & 0.132 & -0.088 & 0.081 & 0.206 \\
\rowcolor{refadj}
METEOR & Extractive & 6 & 0.294 & 0.662 & 0.647 & 0.456 \\
\rowcolor{refadj}
SummaQA & Abstractive & 1 & 0.177 & 0.662 & 0.524 & 0.309 \\
\rowcolor{refadj}
SuPERT & Extractive & 6 & 0.132 & 0.588 & 0.427 & 0.265 \\
\bottomrule
\end{tabular}
\caption{Reference-adjusted: Abstractive or Extractive references.}
\label{tab:traditional_metrics-refadj}
\end{subtable}

\vspace{1em}

\begin{subtable}[t]{\textwidth}
\centering
\begin{tabular}{lllcccc}
\toprule
\textbf{Metric} & \textbf{Ref.} & \textbf{nrefs}
  & \textbf{Coh.} & \textbf{Con.} & \textbf{Flu.} & \textbf{Rel.} \\
\midrule
\rowcolor{reffree}
BLANC & - & 0 & 0.183 & 0.617 & \textbf{0.410} & \textbf{0.383} \\
\rowcolor{reffree}
Length & - & 0 & 0.000 & 0.467 & 0.226 & 0.233 \\
\rowcolor{reffree}
Novel unigram & - & 0 & 0.150 & -0.050 & 0.126 & 0.083 \\
\rowcolor{reffree}
Novel bi-gram & - & 0 & -0.100 & -0.467 & -0.193 & -0.200 \\
\rowcolor{reffree}
Novel tri-gram & - & 0 & -0.100 & -0.533 & -0.259 & -0.267 \\
\rowcolor{reffree}
Repeated unigram & - & 0 & -0.333 & 0.067 & -0.142 & -0.300 \\
\rowcolor{reffree}
Repeated bi-gram & - & 0 & -0.350 & -0.050 & -0.293 & -0.417 \\
\rowcolor{reffree}
Repeated tri-gram & - & 0 & -0.167 & 0.133 & -0.075 & -0.200 \\
\rowcolor{reffree}
Stats-coverage & - & 0 & -0.067 & 0.100 & -0.042 & 0.033 \\
\rowcolor{reffree}
Stats-compression & - & 0 & 0.067 & -0.433 & -0.193 & -0.133 \\
\rowcolor{reffree}
Stats-density & - & 0 & \textbf{0.200} & \textbf{0.633} & 0.360 & 0.367 \\
\bottomrule
\end{tabular}
\caption{Reference-free: nrefs\,=\,0.}
\label{tab:traditional_metrics-reffree}
\end{subtable}
\caption[Traditional (non-LLM) metrics on SummEval.]{%
  \textbf{Traditional (non-LLM) metrics on SummEval.}
  Each metric appears once per setting, using the reference type and nrefs that
  maximise the average Kendall-Tau across the four dimensions.
  \textbf{Bold} indicates the best score within each section per column.
}
\label{tab:traditional_metrics}
\end{table*}


\section{Reference-Adjusted Evaluation on Traditional Metrics}\label{sec:6-chapter-trad-ref-eval}
We conduct a similar analysis to~\Cref{sec:6-chapter-results} with traditional metrics. We compare the use of traditional metrics in the reference-based, reference-adjusted, and reference-free settings. For the reference-based and reference-adjusted settings, we use traditional metrics that support the use of references. For the reference-free setting, we use traditional metrics that do not require a reference. We conduct the same ablations as described in~\Cref{sec:6-chapter-experimental-setup}. We report the results in~\Cref{tab:traditional_metrics}. Interestingly, traditional metrics follow a similar pattern to LLMs-as-judges metrics. While reference-based evaluation generally performs best, reference-adjusted evaluation only moderately decreases the performance from reference-based, and greatly increases the performance over reference-free. While the comparison is not exactly 1:1 (no traditional metrics support both reference-based and reference-free, as LLM judges do), this result indicates that reference-adjusted evaluation may be a viable option for researchers that prefer traditional metrics due to compute or privacy restrictions.  




\section{LLM-as-Judge Results}
\label{sec:6-chapter-results}

We report LLM-as-judge experiments using Prometheus 7b, the best performing model overall, in Table~\ref{tab:prometheus}. We report the results using Mistral 7b and Llama 8b in~\Cref{tab:combined-mistral} and ~\Cref{tab:combined-llama}, respectively. We compare the reference-based, reference-adjusted, and reference-free settings. We additionally conduct a SummEval reproduction study (\Cref{sec:6-chapter-summeval-repro}) and experiment with traditional metrics in all 3 reference settings (\Cref{sec:6-chapter-trad-ref-eval}). We report the best performing traditional metrics for all 3 reference settings for comparison. See~\Cref{tab:traditional_metrics} for the full results on traditional metrics.

The reference-based setting generally outperforms the reference-free, as expected based on past work~\cite{zheng2023judging,kim2023prometheus}. However, we only see the performance advantages of the reference-based setting when multiple references are included. The reference-based setting underperforms the reference-free when only using a single reference. We hypothesize that this drop in performance is due to the model over-indexing on a single reference. Each reference can highlight different parts of the article and the model doesn't get the full picture with a single example. 
This finding is particularly impactful, as collecting gold reference summaries is expensive, and our results suggest that multiple references may be necessary to realize the full benefits of reference-based evaluation, further increasing the cost. 

The reference-adjusted setting leads to performance gains over the reference-free setting. For example, using a single extractive summary with Prometheus raises Consistency correlation from $0.731$ to $0.750$. While using extractive summaries does not consistently outperform the reference-free setting, using abstractive summaries does, and even surpasses the reference-based setting on Coherence and Consistency. We hypothesize that abstractive references provide greater stylistic diversity and are generally higher quality, improving the overall performance of the LLM judge. Notably, the abstractive reference-adjusted setting generally outperforms the reference-based setting when only a single gold reference is available. Since generating multiple synthetic summaries is often cheaper than collecting multiple gold summaries, reference-adjusted evaluation may be a more practical and cost-effective option for low-resource domains.
% -----------------------------------------------------------------------
% TABLE 2: Prometheus Eval
% -----------------------------------------------------------------------
\begin{table*}[h]
\scriptsize
\centering
\setlength{\tabcolsep}{4pt}
\begin{tabular}{lccccc}
\toprule
\textbf{Reference Type} & \textbf{nrefs}
  & \textbf{Coherence} & \textbf{Consistency} & \textbf{Fluency} & \textbf{Relevance} \\
\midrule
\multicolumn{6}{l}{\textit{Reference-based}} \\ \midrule
\rowcolor{refbased}
Gold    & 11 & \textbf{0.726} & \textbf{0.721} & 0.444 & 0.775 \\
\rowcolor{refbased}
Gold    &  6 & 0.701 & 0.642 & \textbf{0.598} & 0.726 \\
\rowcolor{refbased}
Gold    &  1 & 0.548 & 0.682 & 0.296 & \textbf{0.796} \\
\midrule
\multicolumn{6}{l}{\textit{Reference-adjusted}} \\ \midrule
\rowcolor{refadj}
Abstractive & 11 & 0.716 & 0.642 & 0.563 & 0.642 \\
\rowcolor{refadj}
Abstractive &  6 & \textbf{0.750} & 0.509 & \textbf{0.593} & 0.607 \\
\rowcolor{refadj}
Abstractive &  1 & 0.706 & 0.731 & 0.341 & \textbf{0.785} \\
\hdash{1}{6}
\rowcolor{refadj}
Extractive  &  8 & 0.657 & 0.332 & 0.370 & 0.583 \\
\rowcolor{refadj}
Extractive  &  6 & 0.583 & 0.267 & 0.370 & 0.509 \\
\rowcolor{refadj}
Extractive  &  1 & 0.445 & \textbf{0.750} & 0.267 & 0.677 \\
\midrule
\multicolumn{6}{l}{\textit{Reference-free}} \\ \midrule
\rowcolor{reffree}
---         &  0 & \textbf{0.617} & \textbf{0.731} & \textbf{0.444} & \textbf{0.760} \\
\midrule
\multicolumn{5}{l}{\textit{Best traditional (per setting)}} \\ \midrule
\rowcolor{refbased}
Reference-based & - & 0.550\,(ROUGE-L) & 0.700\,(SummaQA) & 0.711\,(ROUGE-4) & 0.700\,(ROUGE-4) \\
\rowcolor{refadj}
Reference-adjusted & - & 0.515\,(ROUGE-L) & 0.691\,(ROUGE-3) & 0.701\,(ROUGE-4) & 0.588\,(ROUGE-4) \\
\rowcolor{reffree}
Reference-free & - & 0.200\,(Stats-density) & 0.633\,(Stats-density) & 0.410\,(BLANC) & 0.383\,(BLANC) \\
\bottomrule
\end{tabular}
\caption[Evaluation results on SummEval for Prometheus across three evaluation settings.]{%
  Evaluation results on SummEval for Prometheus across three evaluation settings.   \colorbox{refbased}{\vphantom{x} Reference-based}: Gold references.
  \colorbox{refadj}{\vphantom{x} Reference-adjusted}: Abstractive or Extractive references.
  \colorbox{reffree}{\vphantom{x} Reference-free}: nrefs\,=\,0.
Best traditional metric per setting, per column, shown for reference.  \textbf{Bold} indicates the best score within each setting per column. Best traditional metrics from~\Cref{tab:traditional_metrics}.
}
\label{tab:prometheus}
\end{table*}

\input{research-chapters/6-chapter/tables/mistral+llama}

We see similar patterns for Mistral and Llama as with Prometheus. In general, we see that while reference-based generally performs the best, reference-adjusted generally outperforms the reference-free setting. For example, the reference-adjusted setting Coherence correlation using Mistral reaches performance parity with the reference-based setting. Using Llama, the Consistency correlation increases from $0.084$ to $0.199$ when going from the reference-free to the reference-adjusted setting. However,  the traditional metrics outperform the Mistral based evaluation. As this is consistent across all the reference settings, this is likely due to the overall performance of Mistral. Prometheus is the best performing model overall and is a finetuned version of Mistral, indicating that evaluation specific finetuning can generally improve weaker models.
